\section{The inference frontier and its precise scope}
\label{sec:intro}
A boundary measurement must carry information to and from the coefficient
being identified. In a damped lattice, that propagation interacts with the
observation window, the distribution of probe durations, the sample budget,
and the elapsed clock. We make those distinctions explicit rather than treating
a conservative inverse modulus as the inference frontier.

There are two positive experiments in this revision. The first keeps exactly
the fixed-window multiscale law of the preceding manuscript. Its sufficient
shrinking-recovery depth improves from $o((\log n)^{1/3})$, and beyond the
$o(\sqrt{\log n})$ improvement proposed in the report, to
$o(\log n/\log\log n)$; the coefficient radius can be polynomial in $n$.
The second changes only the exploration duration law, allowing exponentially
rare long probes with a finite mean. This achieves depth $r\log n$ for an
explicit $r>0$, with a polynomial coefficient radius, while the total clock
is still linear in $n$ with exponentially controlled deviations. The latter
claim is not made under a fixed maximum duration per stage.

\begin{theorem}[Main conclusions, with constants specified below]
\label{thm:main}
Fix the coefficient box of Section~\ref{sec:model}, noise variance $\sigma^2>0$,
a positive exploration floor $\rho$, a bounded initial-state ball, and bounded
predictable baseline forces. Take a product coefficient prior with coordinate
densities bounded below by a common positive constant.

For the fixed-window law $\omega$ of \eqref{eq:compact-law}, let $j_n=J_n+1$.
If $j_n\log(ej_n)=o(\log n)$ and $0<\eta<1$, the exact working posterior
concentrates on $d_{J_n}<n^{-\eta/48}$ uniformly over the declared policy class.
The no-washout clock is bounded by a constant times $n$.

For the finite-mean law $\nu$ of \eqref{eq:long-law}, if $r,s>0$ and
$C_\infty r+8s<1/2$, the same posterior concentrates on
$d_{\lfloor r\log n\rfloor}<n^{-s}$. The clock satisfies the finite-sample
bound \eqref{eq:clock-tail}. In particular, $r=1/(8C_\infty)$ and $s=1/64$
are valid.

For every bounded-input chronological policy with $n$ Gaussian boundary
readouts, the two-point information bound \eqref{eq:global-kl} holds. Thus no
such policy permits uniform fixed-accuracy recovery when $J_n/\log n\to\infty$.
Consequently the attainable depth order in this class is $\log n$; under a
positive minimum inter-readout time it is also $\log t$ in elapsed time.
The constants in the upper and lower directions need not match.
\end{theorem}

The proof is divided into explicit deterministic inversion, chronological
likelihood concentration, and propagation-based information comparison. The
last assertion is a sharpness statement accompanying an attainable positive
rate, not a substitute for a construction.

\subsection{Closest comparisons and retained results}
Du, Nair and Janson~\cite{DNJ2025}, Theorem~1, prove an adaptive Gaussian
linear-regression BvM under information-eigenvalue growth without a deterministic
information limit. Their Corollary~1 also gives a policy-uniform statement under
uniform growth, and Appendix~F discusses obstacles for general parametric
extensions. Accordingly neither random-information Gaussian approximation nor
policy uniformity by itself is claimed as new here. The retained homogeneous
result treats smooth nonlinear mechanical means and arbitrary working priors
on a bounded Hilbert-space initial state; its contrast and nuisance bounds
are verified mechanically. It requires stronger regularity and contrast than
their linear theorem. A posterior Gaussian approximation still does not imply
frequentist coverage of its random center.

Mikhaylov and Mikhaylov~\cite{MM2019,MM2020} reconstruct semi-infinite Jacobi
matrices from discrete-time boundary dynamics. In particular, Theorem~2 and
Proposition~10 of the preprint underlying~\cite{MM2020} connect response data,
positive connecting matrices, and finite coefficient reconstruction. The
response--moment--Jacobi mechanism is therefore classical, not a claimed new
inverse spectral principle. Our continuous-time scalar damping, noisy
chronological experiment, explicit duration allocation, and statistical depth
comparison are different components. The finite-jet condition estimate used
here follows the structural route proposed by the Round~46 report~\cite{R46};
we supply the secant bounds and their all-depth constants rather than citing a
numerical check as proof.

Decay of functions of banded matrices is also established methodology; see
Benzi and Razouk~\cite{BR2007}. We use an elementary two-way Duhamel/path estimate
adapted to the damped block generator and integrate it over all time. The
new inference statement combines that estimate with a finite-mean duration law
and the complete adaptive likelihood, not a claim to invent matrix decay.
Inverse-stability transfer to Bayesian recovery and time-uniform martingale
confidence methods likewise have precedents~\cite{Vollmer,Howard}.

The earlier five-parameter homogeneous quasi-BvM, strong filter jets, memory
Gaussian image, and finite preparation mixture are retained, without altered
hypotheses, in the accompanying source supplement
\texttt{ROUND47\_RETAINED\_RESULTS.tex}. They are not used to prove the new
infinite-dimensional depth theorem. This separates the present central result
from regular finite-dimensional consequences without withdrawing those
consequences. All earlier source files and the controlling report remain in
the repository.

